% =====================================================================
%  Actividad en equipo en clase — Semana 4, Sesión 1
%  "Conteo IFPUG colaborativo — Sistema 'CitasUPA'"
%  Estándares y Métricas para el Desarrollo de Software (E-EMDS-2)
%  UPA — TIID 5° cuatrimestre — Mayo 2026
% =====================================================================
\documentclass[11pt,a4paper]{article}

\usepackage[utf8]{inputenc}
\usepackage[spanish,es-tabla]{babel}
\usepackage[T1]{fontenc}
\usepackage{lmodern}
\usepackage[margin=1.6cm,top=1.5cm,bottom=1.5cm]{geometry}
\usepackage[table]{xcolor}
\usepackage{colortbl}
\usepackage{tabularx}
\usepackage{array}
\usepackage{multirow}
\usepackage{amsmath}
\usepackage{enumitem}
\usepackage{titlesec}
\usepackage{fancyhdr}
\usepackage{tcolorbox}
\tcbuselibrary{skins,breakable}
\usepackage[hidelinks]{hyperref}

\definecolor{accent}{HTML}{0F766E}
\definecolor{accentbg}{HTML}{ECFDF5}
\definecolor{darktext}{HTML}{1F2937}
\definecolor{midtext}{HTML}{4B5563}
\definecolor{lighttext}{HTML}{6B7280}
\definecolor{rowalt}{HTML}{F3F4F6}
\definecolor{redaccent}{HTML}{DC2626}
\definecolor{amber}{HTML}{D97706}
\definecolor{amberbg}{HTML}{FEF3C7}

\renewcommand{\familydefault}{\sfdefault}

\titleformat{\section}
  {\color{darktext}\sffamily\large\bfseries}
  {}{0pt}{}[{\color{accent}\rule{1.2cm}{1.5pt}}]
\titlespacing*{\section}{0pt}{10pt}{4pt}

\titleformat{\subsection}
  {\color{accent}\sffamily\normalsize\bfseries}
  {}{0pt}{}
\titlespacing*{\subsection}{0pt}{6pt}{2pt}

\pagestyle{fancy}
\fancyhf{}
\renewcommand{\headrulewidth}{0pt}
\renewcommand{\footrulewidth}{0pt}
\fancyfoot[L]{\color{lighttext}\footnotesize\itshape EMDS --- Actividad en equipo S4.1: Conteo IFPUG (CitasUPA)}
\fancyfoot[R]{\color{lighttext}\footnotesize\thepage}

\newcommand{\institutionalheader}{%
\noindent
\hfill{\color{lighttext}\small\itshape Estándares y Métricas para el Desarrollo de Software\,---\,E-EMDS-2}\\[0.15em]
{\color{accent}\rule{1.3cm}{2pt}}\\[0.15em]
{\color{darktext}\fontsize{18}{21}\selectfont\bfseries Actividad en equipo\,---\,Semana 4, Sesión 9}\\[0.15em]
{\color{accent}\fontsize{12}{14}\selectfont\bfseries Conteo IFPUG colaborativo \textmd{\color{midtext}--- Sistema ``CitasUPA''}}\\[0.3em]
}

\begin{document}
\thispagestyle{fancy}
\institutionalheader

% ---------- Datos del equipo ----------
\renewcommand{\arraystretch}{1.4}
\noindent
\begin{tabularx}{\linewidth}{|>{\columncolor{accentbg}\bfseries\footnotesize}p{1.8cm}|X|>{\columncolor{accentbg}\bfseries\footnotesize}p{1.2cm}|p{1.6cm}|>{\columncolor{accentbg}\bfseries\footnotesize}p{1.2cm}|p{2.5cm}|}
  \hline
  Equipo \# & & Grupo & & Fecha & \\\hline
\end{tabularx}\\[0.2em]

\noindent
\renewcommand{\arraystretch}{1.5}
\begin{tabularx}{\linewidth}{|>{\columncolor{accentbg}\bfseries\footnotesize}p{2.2cm}|>{\small}p{0.5cm}|>{\small}X|>{\small}p{0.5cm}|>{\small}X|}
  \hline
  \multirow{2}{*}{\bfseries Integrantes} & \textbf{1.} & & \textbf{2.} & \\\cline{2-5}
                                         & \textbf{3.} & & \textbf{4.} & \\\hline
\end{tabularx}

% ---------- Instrucciones ----------
\section*{Instrucciones}

\noindent
\textbf{Modalidad:} equipo designado. \quad
\textbf{Valor:} 10 pts (mismo puntaje para todo el equipo).\\[0.2em]
Apliquen el método \textbf{IFPUG de Puntos de Función} al sistema \emph{CitasUPA} descrito abajo. Trabajen los \textbf{cinco pasos} en orden: identificación, clasificación, UFP, GSC/VAF y AFP. Justifiquen sus decisiones por escrito; el docente revisará la coherencia entre pasos.

\medskip
\noindent\colorbox{amberbg}{\parbox{\dimexpr\linewidth-2\fboxsep}{%
\textbf{\color{amber}Regla de equipo:} no se acepta que una sola persona resuelva la actividad. Cada integrante debe \textbf{firmar} qué paso lideró (parte inferior de la última página).}}

% ---------- Caso ----------
\section*{Caso: sistema \emph{``CitasUPA''}}

\colorbox{accentbg}{\parbox{\dimexpr\linewidth-2\fboxsep}{%
\textbf{Descripción del módulo a estimar.}\\[0.3em]
La Coordinación Estudiantil de la UPA desea un módulo web para que los alumnos agenden citas con tutores académicos.
\begin{itemize}[leftmargin=1.3em,itemsep=0.05em,topsep=0.2em]
  \item \textbf{Funciones que expone al usuario:}
  \begin{enumerate}[leftmargin=1.2em,itemsep=0em,topsep=0.05em,label=(\arabic*)]
    \item Formulario para que el alumno solicite una cita (selecciona tutor, fecha, hora, motivo).
    \item Formulario para que el tutor confirme, reagende o cancele la cita.
    \item Listado de las citas vigentes del alumno (consulta simple sin cálculos).
    \item Listado de las citas vigentes del tutor (consulta simple).
    \item Reporte mensual con totales: citas atendidas, canceladas, reagendadas, no asistencias, y promedio de duración (incluye cálculos).
  \end{enumerate}
  \item \textbf{Datos mantenidos por el sistema:} una tabla \texttt{citas} y una tabla \texttt{usuarios} (alumnos y tutores).
  \item \textbf{Datos de otro sistema:} catálogo \texttt{materias} (proveniente del sistema escolar institucional, solo lectura).
  \item \textbf{Stack tentativo:} Java + Spring Boot + PostgreSQL + frontend React.
  \item \textbf{Tamaño del equipo:} 4 personas. Aplicación intranet (no requiere alta disponibilidad).
\end{itemize}
}}

\clearpage

% ============================================================================
% PASO 1
% ============================================================================
\section*{Paso 1 --- Identificación de componentes \hfill {\normalsize\color{lighttext}\itshape (2 puntos)}}

\noindent
Clasifiquen cada elemento del caso en \textbf{uno} de los cinco tipos IFPUG: ILF, EIF, EI, EO, EQ.

\renewcommand{\arraystretch}{1.5}
\noindent
\begin{tabularx}{\linewidth}{|>{\columncolor{accentbg}\bfseries\small\centering\arraybackslash}p{1.6cm}|>{\small}X|>{\columncolor{accentbg}\bfseries\small\centering\arraybackslash}p{1.4cm}|}
  \hline
  \rowcolor{accent}
  {\color{white}\textbf{\#}} &
  {\color{white}\textbf{Elemento del caso (escriban con sus palabras)}} &
  {\color{white}\textbf{Tipo IFPUG}} \\\hline
  1 & & \\\hline
  2 & & \\\hline
  3 & & \\\hline
  4 & & \\\hline
  5 & & \\\hline
  6 & & \\\hline
  7 & & \\\hline
\end{tabularx}

% ============================================================================
% PASO 2
% ============================================================================
\section*{Paso 2 --- Clasificación por complejidad \hfill {\normalsize\color{lighttext}\itshape (2 puntos)}}

\noindent
Para cada componente del Paso 1, decidan complejidad \textbf{Baja, Media o Alta}. Justifiquen brevemente.

\renewcommand{\arraystretch}{1.5}
\noindent
\begin{tabularx}{\linewidth}{|>{\columncolor{accentbg}\bfseries\small\centering\arraybackslash}p{1.6cm}|>{\columncolor{accentbg}\bfseries\small\centering\arraybackslash}p{1.4cm}|>{\columncolor{accentbg}\bfseries\small\centering\arraybackslash}p{1.8cm}|>{\small}X|}
  \hline
  \rowcolor{accent}
  {\color{white}\textbf{\#}} &
  {\color{white}\textbf{Tipo}} &
  {\color{white}\textbf{Complejidad}} &
  {\color{white}\textbf{Justificación (1 línea)}} \\\hline
  1 & & & \\\hline
  2 & & & \\\hline
  3 & & & \\\hline
  4 & & & \\\hline
  5 & & & \\\hline
  6 & & & \\\hline
  7 & & & \\\hline
\end{tabularx}

\clearpage

% ============================================================================
% PASO 3
% ============================================================================
\section*{Paso 3 --- Cálculo de los UFP \hfill {\normalsize\color{lighttext}\itshape (2 puntos)}}

\noindent
Llenen la matriz con los pesos IFPUG y obtengan los Puntos de Función no Ajustados (UFP).
Pesos IFPUG: \textbf{ILF} 7/10/15; \textbf{EIF} 5/7/10; \textbf{EI} 3/4/6; \textbf{EO} 4/5/7; \textbf{EQ} 3/4/6.

\renewcommand{\arraystretch}{1.5}
\noindent
\begin{tabularx}{\linewidth}{|>{\columncolor{accentbg}\bfseries\small\centering\arraybackslash}p{1.4cm}|X|X|X|X|X|}
  \hline
  \rowcolor{accent}
  {\color{white}\textbf{Tipo}} &
  {\color{white}\centering\textbf{Cant.\ Baja $\times$ peso}} &
  {\color{white}\centering\textbf{Cant.\ Media $\times$ peso}} &
  {\color{white}\centering\textbf{Cant.\ Alta $\times$ peso}} &
  {\color{white}\centering\textbf{Subtotal}} & \\\hline
  ILF & & & & & \\\hline
  EIF & & & & & \\\hline
  EI  & & & & & \\\hline
  EO  & & & & & \\\hline
  EQ  & & & & & \\\hline
  \rowcolor{accentbg}
  \textbf{Total UFP} & \multicolumn{5}{c|}{} \\\hline
\end{tabularx}

% ============================================================================
% PASO 4
% ============================================================================
\section*{Paso 4 --- Evaluación de las 14 GSC y cálculo del VAF \hfill {\normalsize\color{lighttext}\itshape (2 puntos)}}

\noindent
Asignen un valor de \textbf{0 a 5} a cada Característica General del Sistema, según el contexto de \emph{CitasUPA} (intranet, sin alta disponibilidad, equipo pequeño).
Sumen y calculen $\text{VAF} = 0{,}65 + 0{,}01 \cdot \sum \text{GSC}$.

\renewcommand{\arraystretch}{1.4}
\noindent
{\footnotesize
\begin{tabularx}{\linewidth}{|>{\bfseries}p{4.5cm}|>{\centering\arraybackslash}p{1cm}|>{\bfseries}p{4.5cm}|>{\centering\arraybackslash}p{1cm}|}
  \hline
  \rowcolor{accent}
  {\color{white}\textbf{Característica GSC}} & {\color{white}\textbf{0-5}} &
  {\color{white}\textbf{Característica GSC}} & {\color{white}\textbf{0-5}} \\\hline
  1. Comunicaciones de datos     & & 8. Actualización en línea     & \\\hline
  2. Procesamiento distribuido   & & 9. Procesamiento complejo     & \\\hline
  3. Rendimiento                 & & 10. Reusabilidad              & \\\hline
  4. Uso intensivo de config.    & & 11. Facilidad de instalación  & \\\hline
  5. Tasa de transacciones       & & 12. Facilidad de operación    & \\\hline
  6. Entrada de datos en línea   & & 13. Múltiples sitios          & \\\hline
  7. Eficiencia del usuario final& & 14. Facilidad de cambio       & \\\hline
\end{tabularx}}

\medskip
\noindent
\begin{tabularx}{\linewidth}{|>{\columncolor{accentbg}\bfseries\small}p{4cm}|X|>{\columncolor{accentbg}\bfseries\small}p{2cm}|X|}
  \hline
  Suma $\sum \text{GSC}$ & & VAF $=$ & \\\hline
\end{tabularx}

\clearpage

% ============================================================================
% PASO 5
% ============================================================================
\section*{Paso 5 --- AFP y conversión a esfuerzo \hfill {\normalsize\color{lighttext}\itshape (1 punto)}}

\noindent
Calculen los Puntos de Función Ajustados y conviertan a horas de esfuerzo.
Usen un factor de productividad de \textbf{12 horas por AFP} (típico para aplicación web Java/Spring).

\medskip
\renewcommand{\arraystretch}{1.7}
\noindent
\begin{tabularx}{\linewidth}{|>{\columncolor{accentbg}\bfseries}p{4.5cm}|X|}
  \hline
  $\text{AFP} = \text{UFP} \cdot \text{VAF}$ & \\\hline
  Esfuerzo (h) $= \text{AFP} \cdot 12$ & \\\hline
  Equivalencia en meses-persona (1 mes $=$ 160 h) & \\\hline
\end{tabularx}

% ============================================================================
% PASO 6 (reflexión)
% ============================================================================
\section*{Paso 6 --- Análisis de equipo \hfill {\normalsize\color{lighttext}\itshape (1 punto)}}

\noindent
Respondan en una o dos líneas cada una:

\renewcommand{\arraystretch}{1.7}
\noindent
\begin{tabularx}{\linewidth}{|>{\columncolor{accentbg}\bfseries\small}p{6cm}|X|}
  \hline
  ¿Qué componente fue el más difícil de clasificar y por qué? & \\\hline
  ¿Qué GSC les generó más debate? & \\\hline
  Si en vez de intranet fuera una app pública, ¿qué GSC subirían y por qué? & \\\hline
\end{tabularx}

% ============================================================================
% FIRMAS POR PASO
% ============================================================================
\section*{Firma de aporte por integrante}

\noindent\textit{\small\color{lighttext}Cada integrante firma frente a el/los paso(s) que lideró. Si la repartición se ve desbalanceada, se descuenta a todo el equipo.}

\medskip
\renewcommand{\arraystretch}{1.6}
\noindent
\begin{tabularx}{\linewidth}{|>{\columncolor{accentbg}\bfseries\small}p{1.5cm}|X|X|X|X|}
  \hline
  Paso & Integrante 1 & Integrante 2 & Integrante 3 & Integrante 4 \\\hline
  Paso 1 & & & & \\\hline
  Paso 2 & & & & \\\hline
  Paso 3 & & & & \\\hline
  Paso 4 & & & & \\\hline
  Paso 5 & & & & \\\hline
  Paso 6 & & & & \\\hline
\end{tabularx}

% ============================================================================
% VALIDACIÓN
% ============================================================================
\vspace{0.4em}
\section*{Validación}

\renewcommand{\arraystretch}{1.7}
\noindent
\begin{tabularx}{\linewidth}{|>{\columncolor{accentbg}\bfseries\small}p{3.0cm}|X|>{\columncolor{accentbg}\bfseries\small}p{2.4cm}|X|}
  \hline
  Firma del líder & & Firma del docente & \\\hline
  Hora de entrega & & Calificación / 10 & \\\hline
\end{tabularx}\\[0.3em]

\noindent{\footnotesize\color{lighttext}\itshape
Genera 1 firma del rubro \emph{Actividades en clase} (10\,\%) para cada integrante presente. En la revisión grupal compararemos los UFP y AFP obtenidos por cada equipo.}

\end{document}
